% Created 2026-05-26 Tue 09:09
% Intended LaTeX compiler: pdflatex
\documentclass[11pt]{article}
\usepackage[utf8]{inputenc}
\usepackage[T1]{fontenc}
\usepackage{graphicx}
\usepackage{longtable}
\usepackage{wrapfig}
\usepackage{rotating}
\usepackage[normalem]{ulem}
\usepackage{amsmath}
\usepackage{amssymb}
\usepackage{capt-of}
\usepackage{hyperref}
\usepackage{listings}
\usepackage{color}
\usepackage{amsmath}
\usepackage{array}
\usepackage[T1]{fontenc}
\usepackage{natbib}
\author{Hjalte Søberg Mikkelsen}
\date{\today}
\title{}
\begin{document}

\tableofcontents

\section{Abstract}
\label{sec:org7ab9e6a}
Title: A loss-based framework for coherent prediction of competing risks at prespecified time horizons

Keywords: Competing risks, Individualized risk prediction, Loss-based
          estimation, Constrained optimization, Survival analysis, Machine learning

\subsection{Introduction}
\label{sec:orga21d434}
Accurate prediction of disease risk in the presence of competing risks remains a central challenge in clinical epidemiology. We target individualized predictions of the 5-year risk of type 2 diabetes (T2D) and simultaneously the 5-year risk of the competing risk of death. Common estimation methods, such as Fine-Gray regression and cause-specific Cox regression, target only one cause at a time and are not optimized for a given prediction time horizon.

\subsection{Methods}
\label{sec:orgfb22391}
We develop a loss-based estimation framework for the simultaneous prediction of multiple competing risks at a prespecified time horizon. Our approach provides a way to estimate a broad class of existing regression models for competing risks such that the resulting models are optimized for individualized risk prediction. Crucially, the framework enables seamless integration of flexible learner libraries and enables structural constraints during estimation. These may include coherence constraints that guarantee that the predicted risks respect the parameter space and domain-informed shape constraints (e.g., monotonicity in age). Estimation is achieved via constrained optimization of a differentiable loss function, such as the Brier score.

\subsection{Results}
\label{sec:org5e48590}
We evaluate the proposed method in a large Danish registry-based cohort, where it provides individualized 5-year risk predictions for T2D and death without T2D. These predictions may be used for population-level monitoring of type 2 diabetes in Denmark, or to support clinical decision-making, for example in risk stratification of patients. The method demonstrates accurate and well-calibrated estimates of the cause-specific risks.

\subsection{Discussion}
\label{sec:orgeb558e5}
This work introduces a flexible and unified framework for individualized prediction of cause-specific risks at a fixed time horizon. By directly optimizing a loss function at the target time point and enforcing structural constraints, the method yields coherent and clinically interpretable risk estimates. A central strength of the framework is its ability to incorporate flexible learner libraries, allowing it to bridge classical survival models with modern machine learning approaches. At the same time, the inclusion of domain-informed constraints ensures that increased flexibility does not compromise the validity or interpretability of the predictions. Ongoing work explores integration with advanced learners, including boosting and deep neural networks, to further enhance performance.




\section{The problem}
\label{sec:orgc9ee5e9}
Assume we want to predict the risk of diabetes and the risk of death at some timepoint \(\tau\), i.e. we want to predict the risk of each cause in a competing risks scenario. Denote \(T\) as the time between time0 and the date of an event and \(D\in\{0,1,...,K\}\) as the cause of the event, where \(D = 0\) is censoring. Let \(X = (X^1,...,X^p)\) be a p-dimensional vector of baseline covariates, which can both be numeric or categorical.\\[0pt]
Consider a setting where \(D\in\{0,1,2\}\) and \(X\) is 1-dimensional. Define the cumulative hazard for \(D = j\) as
\begin{align*}
\Lambda_j(t|x) = \exp(x\beta_j)\Lambda_{0j}(t)
\end{align*}
where \(\Lambda_{0j}\) is the baseline cumulative hazard for cause j, which is defined as \(\Lambda_{0j}(t) = \int_0^t \lambda_{0j}(s)ds\), where \(\lambda_{0j}\) is the baseline hazard. The baseline hazard is however not estimated when fitting a cause-specific cox model. Now define the survival function using the product integral as
\begin{align*}
S(t|x) = \prod_{s \leq t}\left(1-d\Lambda_1(s|x)-d\Lambda_2(s|x)\right),
\end{align*}
or using the exponential approximation
\begin{align*}
S(t|x) = \exp(-\Lambda_1(t|x)-\Lambda_2(t|x)).
\end{align*}
We can then write the cumulative incidence function for cause \(j\), \(F_j(t|x) = P(T\leq t, D = j|X)\), as
\begin{align*}
F_j(t|x) = \int_0^t S(s-|x)d\Lambda_j(s|x).
\end{align*}
The problem I will look at is that current methods of estimation of the cause-specific risks estimates the cause-specific cumulative incidence functions seperately. The problem with this is that it leads to wrong estimations, which can easily be verified as it can happen that for \(\tau\) and some \(x^*\)
\begin{align}
F_1(\tau|x^*)+F_2(\tau|x^*)>1.
\end{align}
I will in this project develope a method to estimate the risks simultaniously such that the above can't happen.

\subsection{Breslow Estimator}
\label{sec:orgd845e7f}
The problem is most likely to do with the Breslow estimator which is, given n datapoints, one covariate \(x\) and two causes, defined as:
\begin{align*}
\hat{\Lambda}_{01}(t) = \sum_{i = 1}^n\frac{I(\tilde{T}_i\leq t)d_{1i}}{\sum_{j\in R_i}\exp(\hat{\beta}_1x_j)}
\end{align*}
where \(\tilde{T}\) is the observed \(T\), \(d_{1i} = I(D_i = 1)\) and \(R_i = \{j|\tilde{T}_j \geq \tilde{T}_i\}\). 


\section{Constraints}
\label{sec:org1d2aa59}

\subsection{Easy "solution"}
\label{sec:org24437f4}
We can "easily" (maybe hard in practise as the runtime can be long) solve the problem of the cause specific risks summing to more than 1, by just doing normal maximum likelihood under the constraint:
\begin{align}
F_1(\tau|x)+F_2(\tau|x)<1
\end{align}
for all \(x\). This will however only solve that one problem. The bigger problem is that the cause-specific risks are probably also wrongly estimated in other places than just where they sum to more than one - this is just the one place we can verify that they are wrong. Normal partial likelihood already estimates the cause-specific hazards simultationously, but in practise we maximize the partial likelihood seperately for each cause, as it can be shown that the solution is the same as doing it simultaniously. I therefore believe that maximum likelihood is not the way forward. I believe the true problem lies in the Breslow estimator.


\subsection{The easy peasy constraint on S (prodlim)}
\label{sec:orgf26f62b}
Considering the product integral form of the survival function the constraint
\begin{align*}
d\Lambda(t|x)+d\Lambda(t|x) <= 1\text{ , for all t and x}
\end{align*}
will ensure \(S(t|x)>= 0\), where S is defined using the poduct limit estimator. Using the exponential approximaion will automatically ensure that the survival function cant be negative, however this will not ensure that the sum of the cause specific risks are less than 1, which is what this constraint will ensure. This will however only fix the problem where the risks sum to more than one and not the underlying problem.

\section{Different loss functions}
\label{sec:org2bf7682}

\subsection{Likelihood function}
\label{sec:orgcc46972}
Asume we have \(N\) datapoints and \(K\) different causes. The general parametric likelihood function in a competing risks setting can be written as:
\begin{align*}
L = \prod_{i = 1}^N\lambda_{j_i}(t_i|x_i)^{d_i}S(t_i|x_i),
\end{align*}
where \(d_i\) is 1 if an event was observed and 0 if \(i\) is censored, and \(j_i\) is which event occured for person \(i\). Using that the survival function is a product of the individual cause-specific survival functions, \(S(t_i|x_i) = \prod S_j(t_i|x_i)\), we can write the likelihood as:
\begin{align*}
L &= \prod_{i = 1}^N\lambda_{j_i}(t_i|x_i)^{d_i}\prod_{j = 1}^K\exp(-\Lambda_j(t_i|x_i))\\
&= \prod_{i = 1}^N\prod_{j = 1}^K\lambda_{j_i}(t_i|x_i)^{d_{ij}}\exp(-\Lambda_j(t_i|x_i)).
\end{align*}
This means that the likelihood function is a product of K likelihoods, one for each cause. We can then maximize this likelihood by maximizing each cause-specific likelihood function, as long as the parameters are different for each cause.

\subsubsection{Likelihood estimating equations}
\label{sec:org74dda67}
Assume two cause,  only one covariate and constant baseline hazards. I then differentiate the log-likelihood functions with respect to the parameters. The negative log-likelihood function for cause \(j\) is then:
\begin{align*}
-\log(L_j) = \sum_{i = 1}^N-d_{ij}(\beta_jx_i+\log(\lambda_{0j}))+t_i\exp(\beta_jx_i)\lambda_{0j}
\end{align*}
and differentiated
\begin{align*}
\frac{d-\log(L_j)}{d\beta_j} = \sum_{i = 1}^N-d_{ij}x_i+t_ix_i\exp(\beta_jx_i)\lambda_{0j}
\end{align*}
\begin{align*}
\frac{d-\log(L_j)}{d\lambda_{0j}} = \sum_{i = 1}^N-\frac{d_{ij}}{\lambda_{0j}}+t_i\exp(\beta_jx_i)
\end{align*}




\subsection{Brier score}
\label{sec:org42e6b37}
As our target in risk prediction is to minimize the brier-score, it makes sense to fit the models by choosing the parameters that minimize the brier-score. Another good thing about the brier-score is that we use the cumulative incidence function directly in the score, which depends on all the cause specicfic cumulative hazard functions. This may solve the problem without needing a constraint (test).
\\[0pt]
Assuming we have \(N\) datapoints of the form \((\tilde{T},\tilde{D},X)\) where \(\tilde{T}\) is the observed time (time to first observed event), \(\tilde{D}\) is the observed event and \(X\) are the covariate vector. The Brier-Score for predicting cause \(j\) at timepoint \(\tau\) can be written as
\begin{align*}
BS = \frac{1}{N}\sum_{i = 1}^N(F_j(\tau|x_i)-o_i(\tau,j))^2,
\end{align*}
where \(o_i(\tau,j)\) is the true outcome of person \(i\), i.e. \(o_i(\tau,j) = 1(\tilde{D}_i= j,\tilde{T}_i\leq \tau)\). Assume \(K\) different causes and assume no censoring (with censoring just add a 0 state such that we have \(K+1\) states and the setup will be similar to that of the joint survival super learner). We then want to find the parameters that minimize this score i.e.
\begin{align*}
BSS = &\arg\min_{\boldsymbol{\beta}\in\mathbb{R},\boldsymbol{\Lambda_0}\in\mathbb{R}^+}\sum_{j = 1}^K\frac{1}{N}\sum_{i = 1}^N(F_j(\tau|x_i,\boldsymbol{\beta},\boldsymbol{\Lambda_0})-o_i(\tau,j))^2\\
= &\arg\min_{\boldsymbol{\beta}\in\mathbb{R},\boldsymbol{\Lambda_0}\in\mathbb{R}^+}\frac{1}{N}\sum_{j = 1}^K\sum_{i = 1}^N\left(\int_0^{\tau} S(s-|x_i,\boldsymbol{\beta},\boldsymbol{\Lambda_0})d\Lambda_j(s|x_i,\boldsymbol{\beta}_j,\Lambda_{0j})-o_i(\tau,j)\right)^2,
\end{align*}
where \(\boldsymbol{\beta}\) are the beta coeffecients for all causes and \(\boldsymbol{\Lambda_0}\) are the baseline hazards, with j indicating the coeffecients for the specific cause. BSS = Brier-score-sum. This is actually just the Brier-score as it was written by Brier in the orginal paper. 



\subsubsection{Find estimating equations in Parametric Cox with exponential baseline hazard}
\label{sec:orge024b95}
Assume we are in a competing risks setting with two causes and one covariate \(x\), and we are interrested in predicting the risk of experiencing the events before or at timepoint \(\tau\) using a cox model and let \(\phi = \{\beta_1,\beta_2,\Lambda_{01},\Lambda_{02}\}\). First we take the simpler example and use the exponential approximation of the survival function and assume a constant baseline hazard. The hazard for cause \(j\) will the be
\begin{align*}
\lambda_j = \exp(\beta_jx)\lambda_{0j},
\end{align*}
where \(\lambda_{0j}\) is some constant. We are interested in solving this:
\begin{align*}
BSS=\arg\min_{\beta_1,\beta2\in\mathbb{R},\lambda_{01},\lambda_{02}\in\mathbb{R}^+}\frac{1}{N}\sum_{j = 1}^2\sum_{i = 1}^N\left(g_j(\phi,x_i)-o_i(\tau,j)\right)^2.
\end{align*}
First I rewrite \(g_1(\phi) = \int_0^{\tau} S(s-|x,\boldsymbol{\beta},\boldsymbol{\Lambda_0})\lambda_1(s|x,\boldsymbol{\beta}_1,\Lambda_{01})ds\) from the Brier-score written above:
\begin{align*}
g_1(\phi,x)&= \int_0^{\tau}\exp\left[-\exp(\beta_1x)\Lambda_{01}(s-)-\exp(\beta_2x)\Lambda_{02}(s-)\right]\exp(\beta_1x)\lambda_{01}ds\\
&=\exp(\beta_1x)\lambda_{01}\int_0^{\tau}\exp\left[-\exp(\beta_1x)s\lambda_{01}-\exp(\beta_2x)s\lambda_{02}\right]ds\\
&=\exp(\beta_1x)\lambda_{01}\frac{\exp\left[\tau(-\exp(\beta_1x)\lambda_{01}-\exp(\beta_2x)\lambda_{02})\right]-1}{-\exp(\beta_1x)\lambda_{01}-\exp(\beta_2x)\lambda_{02}}\\
&=\exp(\beta_1x)\lambda_{01}\frac{1-\exp\left[\tau(-\exp(\beta_1x)\lambda_{01}-\exp(\beta_2x)\lambda_{02})\right]}{\exp(\beta_1x)\lambda_{01}+\exp(\beta_2x)\lambda_{02}}
\end{align*}
Then we differentiate \(g_1\) with respect to \(\beta_1\), \(\beta_2\), \(\lambda_{01}\) and \(\lambda_{02}\) (to get the result for \(g_2\) just swap all the 1's and 2's, also the ones on the left side of the equals sign):
\begin{align*}
\frac{\partial g_1(\phi,x)}{\partial\beta_1} = &\frac{x \lambda_{01} e^{x \beta_1}}{\left(\lambda_{01} e^{x \beta_1} + \lambda_{02} e^{x \beta_2}\right)^2}\\
&\cdot\left[
\lambda_{02} e^{x \beta_2}\left(1 - e^{-\tau\left(\lambda_{01} e^{x \beta_1} + \lambda_{02} e^{x \beta_2}\right)}\right)
+ \tau \lambda_{01} e^{x \beta_1}\left(\lambda_{01} e^{x \beta_1} + \lambda_{02} e^{x \beta_2}\right)e^{-\tau\left(\lambda_{01} e^{x \beta_1} + \lambda_{02} e^{x \beta_2}\right)}
\right]
\end{align*}
\begin{align*}
\frac{\partial g_1(\phi,x)}{\partial\beta_2} =\frac{xe^{\beta_1x}\lambda_{01}e^{\beta_2x}\lambda_{02}}{(e^{\beta_1x}\lambda_{01}+e^{\beta_2x}\lambda_{02})^2}\left[(1+\tau(e^{\beta_1x}\lambda_{01}+e^{\beta_2x}\lambda_{02}))e^{-(e^{\beta_1x}\lambda_{01}+e^{\beta_2x}\lambda_{02})\tau}-1\right]
\end{align*}
\begin{align*}
&\frac{\partial g_1(\phi,x)}{\partial \lambda_{01}}
=\frac{e^{\beta_1 x}}{\left(e^{\beta_1 x}\lambda_{01}+e^{\beta_2 x}\lambda_{02}\right)^2}\\
&\cdot\left[e^{\beta_2 x}\lambda_{02}\left(1-e^{-\tau\left(e^{\beta_1 x}\lambda_{01}+e^{\beta_2 x}\lambda_{02}\right)}\right)+\tau e^{\beta_1x}\lambda_{01}\left(e^{\beta_1 x}\lambda_{01}+e^{\beta_2 x}\lambda_{02}\right)e^{-\tau\left(e^{\beta_1 x}\lambda_{01}+e^{\beta_2 x}\lambda_{02}\right)}\right]
\end{align*}
\begin{align*}
\frac{\partial g_1(\phi,x)}{\partial\lambda_{02}} = \frac{e^{\beta_1x}\lambda_{01}e^{\beta_2x}}{(e^{\beta_1x}\lambda_{01}+e^{\beta_2x}\lambda_{02})^2}\left[(1+\tau(e^{\beta_1x}\lambda_{01}+e^{\beta_2x}\lambda_{02}))e^{-(e^{\beta_1x}\lambda_{01}+e^{\beta_2x}\lambda_{02})\tau}-1\right]
\end{align*}
When implementing this in R all the things that can be inserted into 1 exp() should be done so terms like exp*exp should be made into 1 exp as otherwise they can explode or be 0*inf in R. Using this we can now differentiate the Brier-Score (BS) with respect to the betas to get the estimating equations:
\begin{align*}
\frac{\partial BSS(\phi,x)}{\partial\beta_1} &=\frac{1}{N}\sum_{j = 1}^K\sum_{i = 1}^N \frac{dg_j(\phi)}{\beta_1}2(g_j(\phi)-o_i(t,j))\\
&= 0
\end{align*}
\begin{align*}
\frac{\partial BSS(\phi,x)}{\partial\beta_2} &=\frac{1}{N}\sum_{j = 1}^K\sum_{i = 1}^N \frac{dg_j(\phi)}{\beta_2}2(g_j(\phi)-o_i(t,j))\\
&= 0
\end{align*}
\begin{align*}
\frac{\partial BSS(\phi,x)}{\partial\lambda_{01}} &=\frac{1}{N}\sum_{j = 1}^K\sum_{i = 1}^N \frac{dg_j(\phi)}{\lambda_{01}}2(g_j(\phi)-o_i(t,j))\\
&= 0
\end{align*}
\begin{align*}
\frac{\partial BSS(\phi,x)}{\partial\lambda_{02}} &=\frac{1}{N}\sum_{j = 1}^K\sum_{i = 1}^N \frac{dg_j(\phi)}{\lambda_{02}}2(g_j(\phi)-o_i(t,j))\\
&= 0
\end{align*}
Note that the \(\frac{1}{N}\) terms can be removed. We also need the hessian:
\begin{align*}
H_{BSS} = 
\begin{bmatrix}
\frac{\partial BSS(\phi,x)}{\partial\beta_1^2} & \frac{\partial BSS(\phi,x)}{\partial\beta_1\partial\beta_2} & \frac{\partial BSS(\phi,x)}{\partial\beta_1\partial\lambda_1} & \frac{\partial BSS(\phi,x)}{\partial\beta_1\partial\lambda_2}\\
\frac{\partial BSS(\phi,x)}{\partial\beta_2\partial\beta_1} & \frac{\partial BSS(\phi,x)}{\partial\beta_2^2} & \frac{\partial BSS(\phi,x)}{\partial\beta_2\partial\lambda_1} & \frac{\partial BSS(\phi,x)}{\partial\beta_2\partial\lambda_2} \\
\frac{\partial BSS(\phi,x)}{\partial\lambda_1\partial\beta_1} & \frac{\partial BSS(\phi,x)}{\partial\lambda_1\partial\beta_2} & \frac{\partial BSS(\phi,x)}{\partial\lambda_1^2} & \frac{\partial BSS(\phi,x)}{\partial\lambda_1\partial\lambda_2} \\
\frac{\partial BSS(\phi,x)}{\partial\lambda_2 \partial\beta_1} & \frac{\partial BSS(\phi,x)}{\partial\lambda_2 \partial\beta_2} & \frac{\partial BSS(\phi,x)}{\partial\lambda_2 \partial\lambda_1} & \frac{\partial BSS(\phi,x)}{\partial\lambda_2^2}
\end{bmatrix}
\end{align*}

\section{Parametric Cox competing risks with binary x (problem doesnt arise)}
\label{sec:org3181b25}
Here I will show that the problem doesnt arise when fitting a parametric cox model using maximum likelihood, where I will use the exponential distribution as the baseline hazard (meaning a constant baseline hazard) (See the Bender paper). Assume a competing risk setting with 2 causes where \(D\in\{0,1,2\}\) is the event indicator. Then let \(x\in\{0,1\}\) be a binary covariate. The cause specific hazard is then
\begin{align*}
\lambda_j(t|x) = \lambda_{0j}\exp(\beta_jx),
\end{align*}
where \(\lambda_{0j}\) is a constant baseline hazard. Note that the hazard doesnt depend on t. The survival function is
\begin{align*}
S(t|x) = \exp\left(-\int_0^s\lambda_1(s|x)+\lambda_2(s|x)ds \right),
\end{align*}
and the density can written as
\begin{align*}
f(t|x) = (\lambda_1(t|x)+\lambda_2(t|x))S(t|x).
\end{align*}

\subsection{Maximum likelihood estimates}
\label{sec:orgd4a1829}
Given n observations, we can write the likelihood as
\begin{align*}
L = \prod_{i = 1}^nf(t_i|x_i)^{\delta_i}S(t_i|x_i)^{1-\delta_i} =\prod_{i = 1}^n (\lambda_1(t_i|x_i)+\lambda_2(t_i|x_i))^{\delta_i}S(t_i|x_i),
\end{align*}
where \(\delta_i = I(D_i \neq 0)\). This notation is a bit wrong as only one of the hazards can contribute to the likelihood at a time if an event was observed, i.e. if person \(i\) had event 1 and was ot censored then the likelihood contribution of person \(i\) would be \(\lambda_1(t_i|x_i)S(t_i|x_i)\). The correct likelihood after using correct notation and rewriting is
\begin{align*}
L = \lambda_{01}^{N_1}\exp(\beta_1\sum_{i\in\delta_{i1}}x_i)\lambda_{02}^{N_2}\exp(\beta_2\sum_{i\in\delta_{i2}}x_i)\exp(-\lambda_{01}\sum_{i = 1}^nt_i\exp(\beta_1x_i))
\exp(-\lambda_{02}\sum_{i = 1}^nt_i\exp(\beta_2x_i)),
\end{align*}
where \(N_j = \sum_{i = 1}^nI(D_i = j)\) and \(\delta_{ij} = \{i|D_i = j\}\). Then let
\begin{align*}
L_j = \lambda_{0j}^{N_j}\exp(\beta_j\sum_{i\in\delta_{ij}}x_i)
\exp(-\lambda_{0j}\sum_{i = 1}^nt_i\exp(\beta_jx_i)),
\end{align*}
and realize that \(\max(L) = \max(L_1)\max(L_2)\). Taking the logorithm of \(L_j\) and differentiating with respect to \(\lambda_{0j}\) and \(\beta_j\)
\begin{align*}
\frac{d\log(L_j)}{d\lambda_{0j}} = \frac{N_j}{\lambda_{0j}}-\sum_{i = 1}^nt_i\exp(\beta_jx_i),
\end{align*}
\begin{align*}
\frac{d\log(L_j)}{d\beta_j} = \sum_{i\in\delta_{i1}}x_i-\lambda_{0j}\sum_{i = 1}^nt_ix_i\exp(\beta_jx_i).
\end{align*}
Setting these to 0 and solving for the two parameters I get
\begin{align*}
\beta_j = \log\left(\frac{\sum_{i\in R'}t_i}{\left(\frac{N_j}{\sum_{i\in\delta_{ij}}x_i}-1 \right)\sum_{i\in R}t_i} \right),
\end{align*}
and
\begin{align*}
\lambda_{0j} = \frac{N_j}{\sum_{i = 1}^n t_i\exp(\beta_jx_i)},
\end{align*}
where \(R = \{i|x_i = 1\}\) and \(R' = \{i|x_i = 0\}\).

\subsection{Seeing if and how F\textsubscript{1} + F\textsubscript{2} >1 (REWRITE AS THE CONSTRAINT CANT BE USED)}
\label{sec:orgd224885}

I will check this by looking at the "easy peasy" constraint, and realize that the only way it doesnt hold is if there exists some T where the following is true:
\begin{align*}
1-d\Lambda_1(T|x)-d\Lambda_2(T|x)<0.
\end{align*}
As the hazards doesnt depend on t \(\Lambda_1(T|x) = T\lambda_1(T|x)\). Inserting this in the above constraint, rewriting a bit and assuming \(x = 0\) we get
\begin{align*}
dT(\lambda_{01}+\lambda_{02})>1,
\end{align*}
where \(dT\) the extremely small interval from the product integral (i.e. infinitly close to 0). Inserting the maximum likelihood estimates of the parameters and rewriting we get
\begin{align*}
\sum_{i\in R'}t_i<dT(N_1+N_2-\sum_{i\in\delta_{i1}}x_i-\sum_{i\in\delta_{i2}}x_i).
\end{align*}
This can ofcourse never be true as \(dT\) is infinitely close to 0, where as the \(t_i\) sum doesnt tend to \(\infty\) as n is a fixed number. This means the sum of the cause specific risks cant exceed 1 in this case.


\section{Parametric Cox competing risks with constant or stepwise constant baseline hazard (poisson regression)}
\label{sec:orged85c94}

\subsection{Constant baseline hazard (exponential distribution)}
\label{sec:org425d4f9}
Consider the sum
\begin{align*}
F_1(\tau|x)+F_2(\tau|x)
\end{align*}
where \(\tau\) is some specific timepoint, i.e. the prediction horizon. Lets write the sum out using the exponential approximation of \(S\) and assuming that the baseline hazards are constant.
\begin{align*}
&F_1(\tau|x)+F_2(\tau|x) = \\
&\int_0^{\tau}\exp(-\Lambda_1(s^-|x)-\Lambda_2(s^-|x))\exp(\beta_1x)\lambda_{01}ds-
\int_0^{\tau}\exp(-\Lambda_1(s^-|x)-\Lambda_2(s^-|x))\exp(\beta_2x)\lambda_{02}ds\\
&= \int_0^{\tau}\exp(-\Lambda_1(s^-|x)-\Lambda_2(s^-|x))ds\left[\exp(\beta_1x)\lambda_{01}+\exp(\beta_2x)\lambda_{02}  \right]\\
&=  \int_0^{\tau}\exp(-\exp(\beta_1x)s^-\lambda_{01}-\exp(\beta_2x)s^-\lambda_{02})ds\left[\exp(\beta_1x)\lambda_{01}+\exp(\beta_2x)\lambda_{02}  \right]\\
&= \frac{\exp(\tau^-[-\exp(\beta_1x)\lambda_{01}-\exp(\beta_2x)\lambda_{02}])-1}{-\exp(\beta_1x)\lambda_{01}-\exp(\beta_2x)\lambda_{02}}[\exp(\beta_1x)\lambda_{01}+\exp(\beta_2x)\lambda_{02}]\\
&= 1-\exp(\tau^-[-\exp(\beta_1x)\lambda_{01}-\exp(\beta_2x)\lambda_{02}])<1
\end{align*}
This can be recognized as \(1-S(\tau|x)\), and means that if the baseline hazard is exponentially distributed then the cause-specific risks can't exceed 1.

\subsection{Stepwise constant baseline hazard (poisson regression)}
\label{sec:orgf634324}

Probably the same as the above, just longer calculations and need to add more notation for the intervals of the baseline hazard.



\section{New way - Brier-score with constraint (should write this into other sections)}
\label{sec:org50f897f}

Assuming we have \(N\) datapoints, the Brier-Score for predicting cause \(j\) at timepoint \(\tau\) can be written as
\begin{align*}
BS_j = \frac{1}{N}\sum_{i = 1}^N(F_j(\tau|x_i)-o_i(\tau,j))^2,
\end{align*}
where \(o_i(\tau,j)\) is the true outcome of person \(i\) (1 if event j is the observed event and has happened at or before timepoint t and 0 otherwise). Furthermore assume only two competing risks (T2D and death). We then want to find the parameters that minimize this score i.e.
\begin{align*}
&\arg\min_{\boldsymbol{\beta}\in\mathbb{R}}\sum_{j = 0}^K\frac{1}{N}\sum_{i = 1}^N(F_j(\tau|x_i,\boldsymbol{\beta})-o_i(\tau,j))^2.
\end{align*}
In our setting \(K = 2\) where \(j = 0\) is the censored state and 1 and 2 are T2D and death. We need to minimize this score under some constraint, for example:
\begin{align*}
\sum_{j = 0}^K F_j(\tau|x,\boldsymbol{\beta})\leq 1 \text{ for all x},
\end{align*}
or
\begin{align*}
\frac{d}{d\beta_{age}}F_j(\tau|x_i,\boldsymbol{\beta})>0
\end{align*}
for some \(j\). In the T2D and death example \(j\in\{1,2\}\), as those risks should increase with age. In practise in may be easier to rewrite the minimization problem to:
\begin{align*}
&\arg\min_{\boldsymbol{\beta}\in\mathbb{R}}\sum_{j = 0}^K\frac{1}{N}\sum_{i = 1}^N(F_j(\tau|x_i,\boldsymbol{\beta})-o_i(\tau,j))^2+ I\left(\frac{d}{d\beta_{age}}F_j(\tau|x_i,\boldsymbol{\beta})\leq0\right),
\end{align*}
i.e. add the constraint to the loss function as a new term which will be 0 if the constraint holds and 1 otherwise which is large as we are using the Brier-score. We should probably multiply the indicator function with \(K+1\), as we sum \(K+1\) Brier-scores. \\[0pt]
If we want methods like Newton-Raphson to work we probably need to use a different method as above, as the loss-function need to be smooth. Say we have the constraint \(g(\beta)\leq 0\) then we can use the loss function:
\begin{align*}
&\arg\min_{\boldsymbol{\beta}\in\mathbb{R}}\sum_{j = 0}^K\frac{1}{N}\sum_{i = 1}^N(F_j(\tau|x_i,\boldsymbol{\beta})-o_i(\tau,j))^2+\lambda\max(0,g(\beta))
\end{align*}
where \(\lambda\) is some penalization parameter. 

\subsection{Easy solution}
\label{sec:org9c988a9}
Fit all learners in the library using normal methods (maximum likelihood, the usual loss functions for the given method, etc.) and then use superlearner methodology to choose the learner that minimizes the above Brier-score loss function and upholds the contraint. This can easily be done with minimal effort by modifying the Score() function in riskRegression. A more advanced version of this is to apply the constraint during the fitting of the individual models. Don't know how hard this is, and by sovling this problem we might aswell extend it to using the Brier-score during the fitting of each learner.
\end{document}
